As is the case with many real-world problems, the quantity of interest changes in time.
For many phenomena some kind of mathematical model for time evolution for the quantity of interest can be formulated that is a more or less accurate representation of the real underlying dynamical process.
The quantity of interest itself can be impossible to measure directly, but typically there can be some other measureable data from which the actual quantity of interest depends on.
As an example, it is difficult to measure cloudiness directly but remote sensing satellites can optically take measurements of the atmosphere from which two-dimensional cloudiness images can be approximately recovered.
Then a model can be formulated to evolve the cloudiness images forward in time.
One such model is the advection-diffusion equation.

\subsubsection{Advection-diffusion equation.}
The advection-diffusion partial differential equation (PDE) is an often used simplification for modeling data where a quantity at least appears to be moving some domain of interest.
The advection-diffusion equation is compromised of the advection term that represents the movement, and diffusion that represent the diffusion of the quantity.

Consider a function $u: \Omega \times I \rightarrow \R$ representing some quantity on spatial domain $\Omega \subset \R^n$, and temporal domain $I = (0, T)$.
The plain advection part is defined as
\begin{equation}\label{eq:advection}
    \frac{\partial u}{\partial t} =  - \mathbf{F} \nabla u,
\end{equation}
where $\mathbf{F} = \mathbf{F}(x,t)$ is the vector-valued advection field that transport $u$ forward in time.
Then the pure diffusion equation would be
\begin{equation}\label{eq:diffusion}
    \frac{\partial u}{\partial t} = \nabla \cdot (D \nabla u)
\end{equation}
with $D = D(x,t)$ being the diffusion coefficient.
In many practical uses $D$ is assumed to be constant in space, and then the diffusion term $\nabla \cdot (D \nabla u)$ simplifies to $D \Delta u$ and diffusion equation~\eqref{eq:diffusion} recovers to the heat equation.

Now by combining~\eqref{eq:advection} and~\eqref{eq:diffusion} gives the advection-diffusion equation
\begin{equation}\label{eq:advectiondiffusion}
    \frac{\partial u}{\partial t} = \nabla \cdot (D \nabla u) - \mathbf{F} \nabla u.
\end{equation}

If the advection field $\mathbf{F}$, and the diffusion coefficient $D$ were known, and the model was perfect $u$ could be solved given some initial condition $u(\cdot, 0) = u_0(\cdot)$, and a boundary condition for the values of $u$ in the Dirichlet type case or the normal derivative of $u$ along the boundary of $\Omega$ in the Neumann type case.
However, measuring advection and diffusion parameters is often not possible and estimates $\overline{\mathbf{F}}$ and $\overline{D}$ need to be extracted e.g.\ from data.
While many uncertainties are involved in the parameter estimation, the model~\eqref{eq:advectiondiffusion} is also not a perfect representation so the approximate model
\begin{equation}\label{eq:advectiondiffusion_approx}
    \frac{\partial u}{\partial t} = \nabla \cdot (\overline{D} \nabla u) - \overline{\mathbf{F}} \nabla u + e
\end{equation}
needs to be considered where $e = e(x,t)$ is the (hopefully small) modeling error.
When taking real measurements, the model also needs to be discretised both in time so that each time step represents one measurement point, and in space so that $\Omega$ becomes a grid in which the measurements are taken at some resolution.

In the cloudiness example the motivation for computing the estimates $\overline{\mathbf{F}}$ and $\overline{D}$ would be that assuming the real parameters change little in time, i.e.\ $\mathbf{F}(x,t_{k}) \approx \mathbf{F}(x,t_{k+1})$ and $D(x,t_{k}) \approx D(x,t_{k+1})$, the parameters can be utilised to transport cloudiness images $u$ to to future and thus producing forecasts.

\subsubsection{Probabilistic state-space models.}
The previously mentioned framework of taking measurements, estimating parameters, and producing forecasts from them can be formalised as the state-space model.
As there are uncertainties related to the measurements, parameter estimation, and the model, instead of studying point estimates of the state it is much more fruitful to consider probability distributions.

Let stochastic process $\{X_k\}_{k=0}^{\infty}$ represent the state of a system which is the quantity of interest, and another stochastic process $\{Y_k\}_{k=0}^{\infty}$ measurements or observations from the system.
Figure~\ref{fig:statespacedependency} presents the dependency scheme of the processes needed to formulate the state-space model.
\begin{figure}[h]
    \begin{alignat*}{5}
        \text{State:}& \quad && X_0 \, \rightarrow \, && X_1 \, \rightarrow \, && X_2 \,\rightarrow \cdots \rightarrow \, && X_k \, \rightarrow \cdots \\
                     & &&  && \downarrow && \downarrow && \downarrow \\
        \text{Observations:}& \quad && && Y_1 && Y_2 && Y_k
    \end{alignat*}
    \caption{\label{fig:statespacedependency}Dependency scheme of the state-space model.}
\end{figure}
For the dependencies to materialise, some assumptions on the processes are needed:
\begin{enumerate}
    \item 
        $\{X_k\}_{k=0}^{\infty}$ satisfies the Markov property with respect to its own history, and history of the observations:
        \begin{equation*}
            \pi (x_{k+1} \, \vert \, x_{0:k}, y_{1:k}) = \pi (x_{k+1} \, \vert \, x_k).
        \end{equation*}
    \item
        $\{Y_k\}_{k=0}^{\infty}$ satisfies the Markov property with respect to the history of the true state:
        \begin{equation*}
            \pi (y_k \, \vert \, x_{0:k}) = \pi (y_k \, \vert \, x_k).
        \end{equation*}
\end{enumerate}
With these assumptions the following conditional probability distributions define the probabilistic state-space model:
\begin{subequations}\label{eq:statespace}
    \begin{align}
        X_0 &\sim \pi (x_0)\\
        X_{k+1} &\sim \pi (x_{k+1} \, \vert \, x_{k}), \quad k = 0,1,2,\dots \label{eq:statespace1}\\
        Y_k &\sim \pi (y_k \, \vert \, x_k), \quad k = 1,2,\dots. \label{eq:statespace2}
    \end{align}
\end{subequations}
Here~\eqref{eq:statespace1} is the model that describes the evolution of state in time, and~\eqref{eq:statespace2} is the observation model that describes the distribution of measurements conditioned on the state.

If functions for the state evolution and observations are known, the same can also be written as 
\begin{subequations}\label{eq:statespacemodel}
    \begin{align}
        X_{k+1} &= M_{k+1}(X_{k}, W_{k+1}), \quad k = 0,1,\dots \label{eq:statespacemodel1}\\
        Y_{k} &= H_{k}(X_{k}, V_{k}), \quad k = 1,2,\dots. \label{eq:statespacemodel2}
    \end{align}
\end{subequations}
Now~\eqref{eq:statespacemodel1} is the state evolution equation where function $M_{k+1}$ gives the next state $X_{k+1}$ given the previous state $X_{k}$ and the state noise $W_{k+1}$
The observation equation~\eqref{eq:statespacemodel2} describes the observations $Y_{k}$ as being the output of the observation function $H_{k}$ with the true state $X_{k}$ and the observation noise $V_{k}$.

With the state-space model~\eqref{eq:statespace} the Bayesian filtering problem can be discussed.
The idea behind Bayesian filtering is to recursively update the predictions and states of the model given latest observations.
The following theorem states the two update equations for the Bayesian filtering scheme.
\begin{theorem}[Bayesian filtering equations]\label{thm:updateformulas}
    Let $\{X_k\}_{k=0}^{\infty}$ be a state process and $\{Y_k\}_{k=0}^{\infty}$ an observation process in a state-space model.
    Then the following recursive equations for Bayesian filtering steps apply:
    \begin{enumerate}
        \item 
            Prediction step:
            \begin{equation}\label{eq:predictionstep}
                \pi (x_{k+1} \, \vert \, y_{1:k}) = \int \pi (x_{k+1} \, \vert \, x_k) \pi (x_k \, \vert \, y_{1:k}) \, d x_k.
            \end{equation}
        \item
            Observation update step:
            \begin{equation}\label{eq:updatestep}
                \pi(x_{k+1} \, \vert \, y_{1:k+1}) = \frac{\pi (y_{k+1} \, \vert \, x_{k+1}) \pi (x_{k+1} | y_{1:k})}{\pi (y_{k+1} \, \vert \, y_{1:k})},
            \end{equation}
            with
            \begin{equation*}
                \pi (y_{k+1} \, \vert \, y_{1:k}) = \int \pi (y_k \, \vert \, x_k) \pi (x_k \, \vert \, y_{1:k-1}) \, d x_{k+1}.
            \end{equation*}
    \end{enumerate}
\end{theorem}
\begin{proof}
    In Appendix~\ref{app:proofs}.
\end{proof}

The Theorem~\ref{thm:updateformulas} provides a framwork for efficiently computing probability distributions conditioned on observations.
The prediction step equation~\ref{eq:predictionstep} being the joint probability distribution of the current and next state conditioned on the latest observations gives a way to get a probabilistic forecast for the future.
Then the observation update equation~\ref{eq:updatestep} is just the Bayes formula where $\pi (x_{k+1} \, \vert \, y_{1:k})$ is interpreted as the prior information of the state for statistical inversion of a new observation.

Bayesian filtering refers to to the Bayesian way of deriving an optimal filtering problem of estimating the state of a nonstationary system through some observations.
The Bayesian filtering theory emerged at the turn of the 1960s with perhaps the most famous example being the works of Rudolf E. Kálmán who published a recursive solution to a linear special case of the problem in~\cite{kalman}.
Meanwhile in the Soviet Union Ruslan Stratonovich studied the more general nonlinear Bayesian filtering problem, see e.g.~\cite{stratonovich2, stratonovich1}.
The simultaneous advances in digital computers allowed the theory to prove itself quickly useful in many applications, an early well-known example of which is the guidance system of the Apollo spacecraft~\cite{mcgee}.

\subsubsection{Kalman filter}
The Kalman filter provides a closed form solution for the update equations of Theorem~\ref{thm:updateformulas} in the special case of a linear state-space model with Gaussian additive noise process.
In this special case the state-space model corresponding to~\eqref{eq:statespacemodel} becomes
\begin{subequations}\label{eq:kalmanstatespace}
    \begin{align}
        X_{k+1} &= M_{k+1}X_{k}  + W_{k+1}, \quad W_{k+1} \sim \mathcal{N}(0, Q_{k+1})\\
        Y_{k} &= H_{k}X_{k} + V_{k}, \quad V_t \sim \mathcal{N}(0, R_k)
    \end{align}
\end{subequations}
with a Gaussian distribution for the initial state $X_0$.

With these assumptions probability densities of the state of the model can be given by just estimates of the mean and covariance.
Now the following theorem gives the filtering step equations of Theorem~\ref{thm:updateformulas}.
\begin{theorem}[Kalman filter equations]
    Assume a linear Gaussian state-space model~\eqref{eq:kalmanstatespace}.
    Let $x_{k\,\vert\,\ell}$ and $P_{k\,\vert\,\ell}$ represent the a posteriori mean and covariance of the state at time $k$ given observations $y_{1:\ell}$.
    Then the filterering equations of Theorem~\ref{thm:updateformulas} take closed form solutions:
    \begin{enumerate}
        \item 
            Prediction step:
            Given
            \begin{equation*}
                \pi (x_k \, \vert \, y_{1:k}) \sim \mathcal{N}(x_{k\,\vert\,k}, P_{k\,\vert\,k}),
            \end{equation*}
            a priori distribution of the next state is
            \begin{equation*}
                \pi (x_{k+1} \, \vert \,  y_{1:k}) \sim \mathcal{N}(\hat{x}_{k+1\,\vert\,k}, \hat{P}_{k+1\,\vert\,k}),
            \end{equation*}
            where parameters of the distribution are given by
            \begin{subequations}\label{eq:kalmanpredictionstep}
                \begin{align}
                    \hat{x}_{k+1\,\vert\,k} &= M_{k+1}x_{k\,\vert\,k}\\
                    \hat{P}_{k+1\,\vert\,k} &= M_{k+1}P_{k\,\vert\,k}M_{k+1}^{\top} + Q_{k+1}.
                \end{align}
            \end{subequations}
        \item
            Observation update step:
            Given
            \begin{equation*}
                 \pi (x_{k+1} \, \vert \,  y_{1:k}) \sim \mathcal{N}(x_{k+1\,\vert\,k}, P_{k+1\,\vert\,k}),
            \end{equation*}
            a posteriori distribution of the next state is
            \begin{equation*}
                 \pi (x_{k+1} \, \vert \,  y_{1:k+1}) \sim \mathcal{N}(x_{k+1\,\vert\,k+1}, P_{k+1\,\vert\,k+1}).
            \end{equation*}
            With 
            \begin{subequations}\label{eq:kalmanupdatestep}
                \begin{align}
                    v_{k+1} &= y_{k+1} - H_{k+1}x_{k+1 \, \vert \, k},\\
                    S_{k+1} &= H_{k+1} \hat{P}_{k+1\,\vert\,k} H_{k+1}^{\top},\\
                    K_{k+1} &= \hat{P}_{k+1\,\vert\,k}  H_{k+1}^{\top} S_{k+1}^{-1},\\
                    x_{k+1\,\vert\,k+1} &= \hat{x}_{k+1\,\vert\,k} + K_{k+1} v_{k+1},\\
                    P_{k+1\,\vert\,k+1} &= (I-K_{k+1}H_{k+1})\hat{P}_{k+1\,\vert\,k}.
                \end{align}
            \end{subequations}
    \end{enumerate}
\end{theorem}
\begin{proof}
    In Appendix~\ref{app:proofs}.
\end{proof}
While the theorem assumes zero means for noises, the Gaussian distributions can easiliy be shifted to accomodate for non-zero means, too.
The Kalman filter relies on the property of Gaussian distributions remaining Gaussian under affine transformations.
However, for nonlinear state-space models one can use linear approximations in order to preserve the convenience of Gaussian densities, a method known as extended Kalman filter.

